\manualtablecaption{Cache-Management Builtins}
\begingroup\footnotesize
\setlength{\tabcolsep}{2pt}
\begin{longtable}{@{}p{1.45in}p{1.3in}p{0.9in}p{1.95in}@{}}
\toprule
\textbf{Name} & \textbf{C interface} & \textbf{Lowering} & \textbf{Availability and effect}\\
\midrule
\endhead
\texttt{flush\_dcache} & \texttt{void(void *,size\_t)} & \texttt{FLSHDCACHE} & Supervisor; applies FLSHDCACHE to every CPUID maintenance block intersecting the range and fully clobbers memory.\\
\texttt{invalidate\_dcache} & \texttt{void(void *,size\_t)} & \texttt{INVDCACHE} & Supervisor; applies INVDCACHE to every CPUID maintenance block intersecting the range and fully clobbers memory.\\
\texttt{invalidate\_icache} & \texttt{void(void *,size\_t)} & \texttt{INVICACHE} & Supervisor; applies local INVICACHE to every CPUID maintenance block intersecting the range, serializes instruction fetch, and fully clobbers memory.\\
\texttt{writeback\_dcache} & \texttt{void(void *,size\_t)} & \texttt{WRBKDCACHE} & Supervisor; applies WRBKDCACHE to every CPUID maintenance block intersecting the range and fully clobbers memory.\\
\texttt{sync\_cache} & \texttt{void(void *,size\_t)} & \texttt{SYNCCACHE} & Supervisor; applies SYNCCACHE to every CPUID maintenance block intersecting the range, serializes instruction fetch, and fully clobbers memory.\\
\bottomrule
\end{longtable}
\endgroup
